\documentclass[11pt]{article}
\usepackage{amsmath,amssymb}
\usepackage[margin=1in]{geometry}
\usepackage{booktabs}
\usepackage{tabularx}
\usepackage{hyperref}
\setlength{\emergencystretch}{2em}

\title{When Are Two Scientific Claims Really the Same?\\
Raw-Evidence Scientific Identity, Plurality, and Integration under Provenance}
\author{ORION Paper III Ultimate-Successor Working Manuscript}
\date{August 2026}

\begin{document}
\maketitle

\begin{abstract}
Scientific knowledge integration fails when lexical, schema, or embedding similarity is mistaken for scientific identity. The bounded ORION Paper III result already shows that an information-equivalent typed product can match ORION, so centralization or graph architecture cannot support the successor headline. P3-U instead targets the acquisition and use of the identity relation from raw heterogeneous evidence. The prospective protocol requires systems to decide whether apparently related scientific objects should be glued, kept distinct, represented as legitimate plurality, marked obstructed, or left unresolved under construct drift, measurement mismatch, context dependence, provenance/stance differences, temporal change, contradictions, and one-to-many mappings. Every decision is source-recoverable and evaluated separately for extraction error, identity-rule error, and insufficient compiled state for the identity responsibility. Strong baselines include modern provenance-aware scientific knowledge-graph extraction, ontology/schema systems, entity matching, RAG synthesis, and donor-complete products receiving equivalent extracted coordinates. Primary outcomes are false merge, false split, plurality preservation, unresolved calibration, provenance fidelity, state/interface-attributable error, expert agreement, and downstream scientific synthesis/QA utility. Superiority requires a prospectively frozen reduction in false scientific merges without an unacceptable false-split/plurality penalty, plus downstream benefit, cross-domain replication, and independent expert/implementation reproduction. Historical ties remain intact; every new false merge or split becomes a responsibility-specific successor experiment rather than an invitation to redefine the gold labels.
\end{abstract}

\section{Residual after donor subtraction}
Scientific literature can already be transformed into provenance-aware knowledge graphs with schema-constrained LLM pipelines \cite{scigraphllm2026}, and end-to-end scientific ontology/KG generation from literature is an active research direction \cite{oarga2026}. Scholarly identity matching itself is also studied as an operational governance problem \cite{ontodup2026}. P3-U therefore claims none of these capabilities as new.

The residual is the protected relation
\[
I(x,y\mid c,e)\in\{\textsc{Glue},\textsc{Distinct},\textsc{Plural},\textsc{Obstruction},\textsc{Unresolved}\},
\]
where $c$ is the scientific context and $e$ is provenance-bound evidence. The question is whether ORION can infer this relation more faithfully from messy science and whether the improvement matters downstream.

\section{Immutable predecessor boundary}
The predecessor's information-equivalent typed product tie is retained as a portability result. It rules out a claim that ORION wins merely because its representation is centralized, graph-shaped, or typed. P3-U must locate any advantage in one of three places:
\begin{enumerate}
    \item better extraction of identity-relevant coordinates from raw evidence;
    \item a better identity/integration rule given the same adequate coordinates;
    \item a better decision about when the available representation is insufficient and must remain unresolved or be reopened.
\end{enumerate}

\section{Identity evidence coordinates}
Candidate identity decisions may use, when present, construct definition, operationalization, measurement procedure, scale/unit, population/context, intervention, causal role, temporal validity, provenance, source stance, uncertainty, and correspondence maps. No coordinate is universally mandatory; the benchmark is designed to reveal which coordinates are responsibility-relevant in each family.

\section{Prospective raw-evidence benchmark}
The corpus is frozen before protected outcomes. It spans several scientific domains and includes double annotation with adjudication. Mandatory families include:

\begin{table}[h]
\centering
\small
\begin{tabularx}{\textwidth}{@{}lXX@{}}
\toprule
Family & Surface trap & Identity obligation \\
\midrule
Construct drift & same term, changed construct & avoid false glue \\
Measurement mismatch & same latent construct, incompatible operationalization & preserve distinction/obstruction \\
Context dependence & same relation under different regimes & condition identity on context \\
Synonymy & different language, same scientific object & avoid false split \\
One-to-many mapping & one coarse object corresponds to several refined objects & preserve plurality \\
Contradictory evidence & sources disagree on compatibility & calibrate unresolved/obstruction \\
Temporal change & identity relation changes after revision/discovery & bind validity in time \\
Provenance/stance & claim and negated/criticized claim share vocabulary & preserve stance \\
State insufficiency & compiled state omits one decisive coordinate & refuse unsupported integration \\
\bottomrule
\end{tabularx}
\caption{Mandatory P3-U identity families.}
\end{table}

Each positive merge has a matched false-merge control; each justified distinction has a matched synonymy/false-split control.

\section{Compiled-view stress test}
The same case is presented through four views:
\begin{enumerate}
    \item raw source evidence;
    \item canonical universal extracted state;
    \item task-compiled state;
    \item responsibility-carrying state that declares which identity obligation it is sufficient to answer.
\end{enumerate}
If two views are certified semantics-preserving and sufficient for the identity responsibility, the promoted decision must be invariant. If a compact state is adequate for QA or prediction but omits a coordinate needed for identity, the correct terminal is not forced glue or split. It is state insufficiency / \textsc{Unresolved}, and the raw source must be reopenable.

\section{Source recoverability}
Every promoted \textsc{Glue}, \textsc{Distinct}, \textsc{Plural}, or \textsc{Obstruction} decision must cite the source spans/records that carry its load-bearing coordinates. A generated ontology label or embedding score cannot self-authorize identity. Representation/schema changes invalidate or transport existing identity certificates through an explicit bridge rather than silently preserving them.

\section{Baselines}
All systems receive matched model, evidence access, and resource budgets. Required comparators include:
\begin{enumerate}
    \item lexical/embedding entity matching;
    \item schema/ontology alignment;
    \item provenance-aware scientific KG extraction in the style of modern constrained pipelines \cite{scigraphllm2026};
    \item LLM raw-text pair judge;
    \item RAG synthesis with entity canonicalization;
    \item donor-complete product receiving the same adequate extracted coordinates as ORION;
    \item raw-text end-to-end arms for each strong system when runnable.
\end{enumerate}
The coordinate-equalized arm isolates the identity rule; raw-text arms measure full-system value.

\section{Primary hypotheses}
\paragraph{H1, false-merge superiority.}
ORION reduces false scientific merges versus the strongest runnable comparator by a frozen prospective margin.

\paragraph{H2, non-compensatory false-split/plurality guard.}
The H1 gain is invalid if false splits or loss of legitimate plurality exceeds frozen non-inferiority guards.

\paragraph{H3, downstream value.}
The improved identity portrait yields higher protected scientific QA/synthesis utility under matched downstream models and evidence.

\paragraph{H4, state-responsibility calibration.}
ORION reduces identity errors attributable to insufficient compiled state by reopening or abstaining when the view is not sufficient.

\paragraph{H5, transfer.}
The H1--H4 pattern replicates across domains and an independently annotated/implemented evaluation.

\section{Statistics and annotation authority}
The scientific item/pair, not repeated model samples, is the default independent unit. Before protected outcomes, freeze sampling, family weights, adjudication procedure, agreement statistic, superiority/non-inferiority margins, multiplicity handling, confidence-interval method, and treatment of unresolved cases.

Report false-merge and false-split effects separately with uncertainty. Do not collapse them into a single F1 score that can hide asymmetric scientific harm. Report expert agreement and disagreement strata; adjudicated labels do not erase genuine ambiguity.

\section{Adversarial controls}
Mandatory attacks include:
\begin{enumerate}
    \item identical terminology with scientifically different constructs;
    \item different terminology for the same construct;
    \item schema-compatible objects with incompatible measurement semantics;
    \item same typed graph but different provenance/stance;
    \item globally information-equivalent states where one interface hides the decisive identity coordinate;
    \item compact state adequate for QA but inadequate for integration;
    \item source-round-trip failure after a promoted merge;
    \item LLM judge agreement caused by shared prompt/template leakage.
\end{enumerate}

\section{Negative-result research recursion}
Each false merge/split/plurality loss is attributed to one primary candidate cause:
\[
\{\text{extraction},\text{identity rule},\text{state/interface sufficiency},\text{provenance},\text{temporal validity},\text{annotation ambiguity},\text{benchmark leakage}\}.
\]
A child experiment changes one suspected cause while holding the others fixed. For example, an extraction error is re-tested with gold coordinates before changing the identity rule; a state-sufficiency error is re-tested by restoring the omitted coordinate before inventing a new ontology relation.

ORION issue \#651 is the authority-bearing P3-U recovery programme. New child issues are opened only for a distinct discovered coordinate/mechanic with a frozen discriminator. Gold relabeling after viewing system errors, post-hoc exclusion of hard pairs, and baseline weakening are prohibited.

\section{Claim ladder and current status}
The manuscript progresses only through: protocol frozen; raw-identity mechanism supported; protected false-merge superiority with guards; downstream utility supported; cross-domain independently replicated bounded superiority. A negative, tie, or \textsc{CannotCheck} terminal remains visible.

This TeX file is prospective. It does not alter the predecessor P3 result and asserts no raw-literature superiority before protected execution.

\begin{thebibliography}{9}
\bibitem{scigraphllm2026}
``SciGraph-LLM: Automatic Knowledge Graph Construction from Scientific Papers,'' Proceedings of WSDM 2026, doi:10.1145/3779211.3793169.

\bibitem{oarga2026}
A. Oarga, M. Hart, A. M. Bran, M. Lederbauer, and P. Schwaller, ``Scientific knowledge graph and ontology generation using open large language models,'' \emph{Digital Discovery}, vol. 5, p. 1269, 2026, doi:10.1039/D5DD00275C.

\bibitem{ontodup2026}
``OntoDup: Governance-Aware Entity Matching for Scholarly Knowledge Graph Deduplication,'' \emph{Information}, vol. 17, no. 4, 325, 2026.
\end{thebibliography}

\end{document}
